\section{A witnessed-choice sentence on all inputs}\label{sec:sentence}

We construct one sentence over the eight binary relations of
Section~\ref{str:structures}. Its computation first orders the graph of
blocks using a tuple of four vertex blocks. It then selects one edge atom
on each edge outside a spanning tree and tests whether those selections
extend to compatible configurations. The choices are guarded by an
explicit test for suitable witnessing automorphisms. Consequently the
construction is meaningful on arbitrary finite inputs, including inputs
whose block graph is not a box.

The state is a set of selected atoms, and a real step adds one atom from
the next unselected edge outside the tree. We permit that step only when
an explicitly constructed family of full-input automorphisms fixes the
selected atoms individually and acts transitively on a nonempty candidate set.
Otherwise a pure dummy leaves the state unchanged. Since states only
grow, pointwise fixation of the current selected atoms fixes every
earlier state as well. This is the invariant that makes the guarded
computation legal on all inputs. On the boxes, the additional work in
Section~\ref{sec:grid-behavior} proves that the guard always permits a
real choice until every non-tree edge has been processed.

\subsection{Bounded notation, blocks, and distances}

Let \(A\) be the universe of an arbitrary finite \(\tau\)-structure, and
write \(n=|A|\). All the definitions below abbreviate BGS terms and
formulas, except for the one explicitly displayed WSC operator. We use
Kuratowski ordered pairs
\[
 \langle a,b\rangle=\{\{a\},\{a,b\}\}
\]
and nested pairs for tuples of any fixed length. A relation is represented
by its set of ordered pairs. Fixed products are formed by nested
comprehensions and finite unions. Bounded quantifiers are abbreviations:
for example, \(\exists z\in t\,\theta(z)\) tests whether
\(\{\emptyset:z\in t,\theta(z)\}\) is nonempty. Membership, intersection,
difference, singletons, and finite unions are therefore available from
\(\Pair,\Union,\Unique\), and comprehension. The predicate that an object
is a set is its nonmembership in \(\Atoms\).

Natural numbers are von Neumann ordinals. Their order is membership, and
their successor is \(i\cup\{i\}\). In particular, the bounded index set
\[
 I_n=\Card(\Atoms)\cup\{\Card(\Atoms)\}=\{0,\ldots,n\}
\]
is a term. Unique lookups are implemented by \(\Unique\); an unsuccessful
lookup returns \(\emptyset\). A conditional term is \(\Unique\) applied
to the union of the two appropriately guarded singleton comprehensions.
Tuple components can equivalently be accessed by bounded tests on their
defining product, so no unbounded search is implicit in this notation.
Every fixed list of conditions is expanded as a finite conjunction or
disjunction.

Define the two sets of atoms
\[
 Y=\{s\in\Atoms:\Ed(s,s)\},\qquad
 X=\{t\in\Atoms:\Cf(t,t)\},
\]
and the two families of rows
\[
 \mathcal E=\bigl\{\{s\in Y:\EB(a,s)\}:a\in Y\bigr\},\qquad
 \mathcal V=\bigl\{\{t\in X:\VB(a,t)\}:a\in X\bigr\}.
\]
These terms are defined before any structural assumptions are tested.
Each family has at most \(n\) members and every row has at most \(n\)
atoms. When the checks below succeed, their members are precisely the
nonempty edge blocks and vertex blocks. Define
\[
 v\sim e
 \quad\Longleftrightarrow\quad
 \exists t\in v\,\exists s\in e\;\Inc(t,s)
 \qquad(v\in\mathcal V,\ e\in\mathcal E).
\]
For \(J\subseteq\mathcal E\), write \(u\sim_Jw\) when \(u,w\) are
distinct vertices and some \(e\in J\) is incident with both.

Here is a bounded term used for every connectivity test in the
construction. Starting with \(R=\emptyset\), iterate the inflationary
update
\begin{align}
 R\longmapsto R
 &\cup\{\langle v,v,0\rangle:v\in\mathcal V\}\notag\\
 &\cup\bigl\{\langle v,w,i+1\rangle:
       v,w\in\mathcal V,\ i,i+1\in I_n,\ 
       \exists u\in\mathcal V\,
       [\langle v,u,i\rangle\in R\ \land\
        (u=w\ \lor\ u\sim_Jw)]\bigr\}.
 \label{sen:distance-iteration}
\end{align}
Denote its fixed point by \(R_J\), and put
\[
 \operatorname{conn}_J(v,w)
 \quad\Longleftrightarrow\quad
 \langle v,w,n\rangle\in R_J.
\]
Every iterate lies in the fixed triple domain
\(\mathcal V^2\times I_n\). Induction on \(i\) shows that its fixed point
contains \(\langle v,w,i\rangle\) exactly when there is a \(J\)-walk from
\(v\) to \(w\) of length at most \(i\): the option \(u=w\) pads a shorter
walk. Since \(|\mathcal V|\leq n\), the predicate
\(\operatorname{conn}_J\) expresses reachability for the adjacency relation \(\sim_J\). Define \(d(v,w)\) to be the least \(i\in I_n\) with
\(\langle v,w,i\rangle\in R_{\mathcal E}\), using \(0\) if there is no
such \(i\).

Let \(\operatorname{Check}(A)\) be the conjunction of the following
bounded conditions.
\begin{enumerate}[label=(\roman*)]
 \item The relations \(\Ed,\Cf\) are exactly the diagonals of a
       partition \(Y\mathbin{\dot\cup}X=A\).
 \item The relation \(\EB\) is an equivalence relation on \(Y\) and
       is empty outside \(Y^2\); likewise \(\VB\) is an equivalence
       relation on \(X\) and is empty outside \(X^2\).
       Also \(\Inc\subseteq X\times Y\).
 \item The family \(\mathcal V\) is nonempty. Every \(e\in\mathcal E\)
       has exactly two distinct incident vertices, and different
       edges have different unordered pairs of endpoints.
       Every two vertices satisfy
       \(\operatorname{conn}_{\mathcal E}(v,w)\).
 \item If \(v\sim e\), then every \(t\in v\) has exactly one
       \(\Inc\)-neighbor in \(e\).
\end{enumerate}
When this check succeeds, let \(G\) denote the resulting simple connected
graph on \(\mathcal V\), with edge set \(\mathcal E\). No condition is
imposed on the three input relations \(\Zrel_\delta\); they will all be
included in the automorphism test.

\subsection{Ordering the block graph and choosing a tree}

A tuple \(\alpha=(a_1,a_2,a_3,a_4)\in\mathcal V^4\) is
\emph{resolving} if the map
\[
 v\longmapsto
 (d(a_1,v),d(a_2,v),d(a_3,v),d(a_4,v))
\]
is injective on \(\mathcal V\). Let \(P\) be the set of codes of all
resolving tuples when \(\operatorname{Check}(A)\) holds, and let
\(P=\emptyset\) otherwise. Repetitions among the four entries are
permitted. For \(\alpha\in P\), lexicographic comparison of the four
distance coordinates defines a linear order \(<_\alpha\) on vertices.
Order edges by the lexicographic order of their endpoint pairs, placing
the smaller endpoint first.

Let \(r\) be the least vertex. At every vertex \(v\ne r\), choose as
parent the least neighbor \(w\) satisfying
\(d(r,w)+1=d(r,v)\), and let \(T=T(\alpha)\) be the set of these parent
edges. Here \emph{least} always means a bounded test that no smaller
eligible member of \(\mathcal V\) or \(\mathcal E\) exists. Connectedness
ensures that a parent exists. Parent distances strictly decrease, so
every parent chain reaches \(r\); consequently \(T\) is a spanning tree.
Set \(T=\emptyset\) for \(\alpha\notin P\).
All later tree constructions are guarded by \(\alpha\in P\).

\subsection{Comparisons on four-cycles and the next choice}

We next define comparisons between atoms on the sides of a four-cycle.
They determine which atoms on the next edge are compatible with the
previously selected atoms. On the boxes, one such comparison detects
equality of the vector attached to that face. Comparing two atoms on the
same edge over all its faces therefore detects equality of their complete
vector coordinates, leaving their scalar coordinates unrestricted.
Lemma~\ref{grid:matching} proves these assertions. The definitions here
use only incidence and apply to arbitrary checked inputs.
Let \(\mathcal F\) consist of all
codes
\[
 f=(v_0,v_1,v_2,v_3;e_0,e_1,e_2,e_3)
 \in\mathcal V^4\times\mathcal E^4
\]
whose four vertices and four edges are separately distinct, with \(e_i\)
joining \(v_i\) to \(v_{i+1}\), where subscripts are taken modulo \(4\).
Every starting point and direction is included.

For adjacent positions \(i,j\), define
\(\operatorname{adj}_f(i,s;j,s')\) by requiring \(s\in e_i\),
\(s'\in e_j\), and a configuration in their shared vertex block incident
with both atoms. Explicitly, the shared vertex is \(v_{i+1}\) when
\(j=i+1\), and \(v_i\) when \(j=i-1\), and the last condition is
\[
 \exists t\in v\;[\Inc(t,s)\land\Inc(t,s')].
\]
For arbitrary positions, define \(\operatorname{match}_f(s,s')\) as the
disjunction over the sixteen pairs \(i,j\in\{0,1,2,3\}\) of \(s\in e_i\),
\(s'\in e_j\), and the following condition:
\begin{itemize}
 \item if \(j=i\pm1\), use
       \(\operatorname{adj}_f(i,s;j,s')\);
 \item if \(j=i\) or \(j=i+2\), use
       \[
       \exists q\in e_{i+1}\;
       [\operatorname{adj}_f(i,s;i+1,q)\land
        \operatorname{adj}_f(j,s';i+1,q)].
       \]
\end{itemize}
The two adjacency tests in the second clause may have different
configuration witnesses. For \(s,s'\in h\in\mathcal E\), put
\begin{equation}
 \operatorname{baseEq}_h(s,s')
 \quad\Longleftrightarrow\quad
 \forall f\in\mathcal F\;
 [h\text{ is a side of }f\ \Longrightarrow\
       \operatorname{match}_f(s,s')].
 \label{sen:base-equality}
\end{equation}

A state \(x\) of the choice computation will be a set of selected edge
atoms. Let \(E_{\rm next}(\alpha,x)\) be the singleton consisting of the
least edge in
\[
 \{e\in\mathcal E\setminus T:x\cap e=\emptyset\},
\]
or the empty set if this set is empty or \(\alpha\notin P\).
Define \(D(\alpha,x)\) by taking, for \(e\in E_{\rm next}(\alpha,x)\),
all \(s\in e\) such that
\begin{equation}
 \operatorname{match}_f(s,s')
 \quad\text{whenever }\
 f\in\mathcal F,\ e\text{ is a side of }f,\ s'\in x,
 \text{ and }s'\text{ lies on a side of }f.
 \label{sen:candidates}
\end{equation}
If there is no next edge, then \(D(\alpha,x)=\emptyset\).
All side tests are finite disjunctions over the four positions of the
tuple \(f\). Thus \(D\) is a bounded comprehension even when the input
fails the structural check.

\subsection{Explicit pools of relations on atoms}

To witness a choice we shall test compositions of two kinds of relations.
Their roles can be seen on a box. To send one candidate edge atom to
another, first translate their differing face vectors. At most four faces
are involved, so this changes only a bounded number of atoms. It may also
change the scalar coordinate; a circulation along the new edge and its
return path in \(T\) corrects the remaining scalar difference. This cycle
contains no previously selected non-tree edge. Proposition~\ref{grid:transitivity}
proves that the two transformations fix all previous selections and
preserve the full structure.

The first pool below contains every permutation moving at most a fixed
number of atoms. The second describes the tree-cycle corrections through
incidence and the relations \(\Zrel_\delta\). These definitions require
neither coordinates nor a promise that the input is a box. They may
produce relations that are not functions; we retain a composition only
after checking that it is an automorphism of the entire input fixing
the selected atoms individually.

\paragraph{Local relations.}
Write
\[
 I_A=\{\langle a,a\rangle:a\in A\},\qquad
 B=16(3^9+3^{29}).
\]
For \(\bar p,\bar q\in A^B\), define the relation
\begin{equation}
 \ell_{\bar p,\bar q}
 =\{\langle p_i,q_i\rangle:1\leq i\leq B\}
   \cup
   \{\langle a,a\rangle:a\in A,\ a\ne p_i
                                      \text{ for every }i\}.
 \label{sen:local-relation}
\end{equation}
Set
\[
 \mathcal L=\{I_A\}\cup
       \{\ell_{\bar p,\bar q}:\bar p,\bar q\in A^B\}.
\]
The indices in Equation~\eqref{sen:local-relation} are a fixed finite
list, not an input-dependent iteration. If a nonidentity permutation
moves at most \(B\) atoms, listing its moved atoms and their images and
padding both lists by repetitions puts its graph in \(\mathcal L\).
Identity was included separately, also covering \(A=\emptyset\).
Other members of \(\mathcal L\) need not be functions.

\paragraph{Cycle relations.}
We now define the second pool \(\mathcal C(\alpha,x)\). It contains
\(I_A\). For \(\alpha\in P\), add a relation for every
\[
 e\in E_{\rm next}(\alpha,x),\quad
 (U,V)\text{ an ordering of the endpoints of }e,\quad
 o\in\mathcal V.
\]
For these indices, the directed traversal consists of \(e\) from \(U\)
to \(V\) and the tree path from \(V\) to \(U\). We specify it without
listing a path as an unbounded tuple. Its set of directed edges is
\begin{align}
 \Gamma_{e,U,V}={}&\{\langle e,U,V\rangle\}\notag\\
 &{}\cup
 \{\langle h,r,r'\rangle:
 h\in T,\ r,r'\text{ are the distinct endpoints of }h,\notag\\
 &\hspace{40mm}
 \neg\operatorname{conn}_{T\setminus\{h\}}(V,U),\
 \operatorname{conn}_{T\setminus\{h\}}(V,r),\
 \operatorname{conn}_{T\setminus\{h\}}(U,r')\}.
 \label{sen:cycle-traversal}
\end{align}
Deleting a tree edge gives two components. It belongs to the return
path precisely when they separate \(V\) and \(U\), and then the
conditions in Equation~\eqref{sen:cycle-traversal} determine its
direction uniquely. Write \(E(\Gamma)\) for the set of edge blocks
occurring in this traversal, and for
\(\langle h,r,r'\rangle\in\Gamma\) set
\[
 \delta_{\Gamma,o}(h)=
 \begin{cases}
  1,&d(o,r)<d(o,r'),\\
  2,&\text{otherwise}.
 \end{cases}
\]
The values \(1,2\) here select the corresponding relation symbols
\(\Zrel_1,\Zrel_2\). In particular, ties in distance are assigned
the value \(2\); no bipartiteness assumption is needed in this definition.
On the box structures of Section~\ref{sec:grid-behavior}, the block at
the coordinate origin is an index \(o\) for which these labels give the
signs of traversal in the fixed edge orientations. Enumerating all
\(o\) includes this index without selecting it;
Proposition~\ref{grid:transitivity} verifies the resulting circulation.

The edge part of the indexed relation is
\begin{align}
 c_Y={}&
 \{\langle s,s\rangle:
       h\in\mathcal E\setminus E(\Gamma),\ s\in h\}\notag\\
 &{}\cup
 \{\langle s,s'\rangle:
       h\in E(\Gamma),\ s,s'\in h,\
       \operatorname{baseEq}_h(s,s')\land
       \Zrel_{\delta_{\Gamma,o}(h)}(s,s')\}.
 \label{sen:cycle-edge-relation}
\end{align}
Its configuration part is
\begin{align}
 c_X=\{\langle t,t'\rangle:\;&v\in\mathcal V,\ t,t'\in v,\notag\\
 &\forall h\in\mathcal E\
 [v\sim h\ \Longrightarrow\
   \exists s,s'\in h\,
    (\Inc(t,s)\land\Inc(t',s')\land
                           \langle s,s'\rangle\in c_Y)]\}.
 \label{sen:cycle-configuration-relation}
\end{align}
Put \(c_{e,U,V,o}=c_Y\cup c_X\), and collect all these relations together
with identity into \(\mathcal C(\alpha,x)\). If \(\alpha\notin P\),
take \(\mathcal C(\alpha,x)=\{I_A\}\). The displayed condition involving
\(\Zrel_{\delta_{\Gamma,o}(h)}\) is a two-case finite disjunction, so
it does not require a relation symbol with a variable index.

\paragraph{Composition and the automorphism filter.}
The two raw pools are now defined. We retain their compositions that
preserve the full input and fix the selected atoms.
For relations \(\ell,c\subseteq A^2\), their composition in this order is
\[
 c\circ\ell
 =\{\langle a,b\rangle:a,b\in A,\
       \exists u\in A\,
       [\langle a,u\rangle\in\ell\land\langle u,b\rangle\in c]\}.
\]
Let \(\operatorname{Filt}(\alpha,x,g)\) require the following.
\begin{enumerate}[label=(\roman*)]
 \item The relation \(g\) is the graph of a permutation of \(A\):
       every atom has exactly one image and exactly one preimage.
 \item For every one of the eight symbols \(R\in\tau\) and every
       \(a,a',b,b'\in A\) with
       \(\langle a,b\rangle,\langle a',b'\rangle\in g\),
       \[
                     R(a,a')\ \longleftrightarrow\ R(b,b').
       \]
 \item Each member of \(\mathcal E\cup\mathcal V\) is fixed setwise,
       and \(\langle s,s\rangle\in g\) for every \(s\in x\cap A\).
\end{enumerate}
Each condition is bounded. For example, setwise preservation of a block
\(u\) is tested by requiring that the image of every \(a\in u\) belongs
to \(u\); bijectivity and finiteness then give equality of the image with
\(u\). Define
\begin{equation}
 M(\alpha,x)=
 \{I_A\}\cup
 \{c\circ\ell:\ell\in\mathcal L,\ c\in\mathcal C(\alpha,x),\
       \operatorname{Filt}(\alpha,x,c\circ\ell)\}
 \qquad(\alpha\in P),
 \label{sen:witness-pool}
\end{equation}
and set \(M(\alpha,x)=\{I_A\}\) otherwise.
The local relation is applied first. Only the composition is filtered;
neither raw factor is assumed to be a permutation.

Every member of \(M(\alpha,x)\), when \(\alpha\in P\), is thus an
automorphism of the full \(\tau\)-structure. It fixes \(\alpha\), since
it fixes all four of its block entries, and it fixes every atom in the
state \(x\) individually whenever \(x\subseteq Y\). This last condition
will ensure that witnesses respect the entire computation history.

\subsection{The guarded choice and witness terms}

Define the bounded formula
\[
 \operatorname{Trans}(\alpha,x)
 \quad\Longleftrightarrow\quad
 \forall s,s'\in D(\alpha,x)\;
   \exists g\in M(\alpha,x)\;
                    \langle s,s'\rangle\in g,
\]
and the choice term
\begin{equation}
 s_{\rm choice}(\alpha,x)=
 \begin{cases}
 D(\alpha,x),&
   D(\alpha,x)\ne\emptyset\ \land\operatorname{Trans}(\alpha,x),\\
 \{\emptyset\},&\text{otherwise}.
 \end{cases}
 \label{sen:guarded-choice}
\end{equation}
The singleton \(\{\emptyset\}\) supplies a pure dummy choice. The step and
witness terms are
\begin{equation}
 \begin{split}
 s_{\rm step}(\alpha,x,y)
      &=x\cup\bigl(\{y\}\cap Y\bigr),\\
 s_{\rm wit}(\alpha,z,x)&=M(\alpha,x).
 \end{split}
 \label{sen:step-witness}
\end{equation}
In the second line \(z\) denotes the resulting fixed point and \(x\)
denotes the intermediate state being certified. These are precisely
the two slots, in this order, of the witnessing term in
\cite[Section~3.1]{ls2023}; the final-state slot is unused.
Equivalently the syntactic term supplied in the slots \((\alpha,x,y)\)
is \(M(\alpha,y)\). Unused formal variables may be inserted through
tautological self-equalities if required. Under the member-choice
convention, every real choice \(y\) is an atom of \(Y\), so the first
line adds exactly that atom. The pure dummy \(\emptyset\) is not in
\(Y\) and causes no change.

\subsection{Testing the terminal selection}

It remains to define the output formula. Say
\(\operatorname{Complete}(\alpha,x)\) holds if \(\alpha\in P\), \(x\)
is a set contained in \(Y\), and
\[
 |x\cap e|=1\quad(e\in\mathcal E\setminus T),\qquad
 x\cap h=\emptyset\quad(h\in T).
\]
These cardinality tests use \(\Card\). When
\(\operatorname{Complete}(\alpha,x)\) holds, the selections on non-tree
edges impose restrictions on configurations within each single block
\(v\). Once these are imposed, the remaining compatibility
conditions connect configurations only across the edges of \(T\).
We test them by successively deleting configurations that have no
compatible survivor across an adjacent tree edge.
Specifically, for \(t\in v\in\mathcal V\), say that \(t\) is
\emph{allowed} if
\begin{equation}
 \forall h\in\mathcal E\setminus T\
 [v\sim h\ \Longrightarrow\
           \exists s\in x\cap h\;\Inc(t,s)].
 \label{sen:allowed}
\end{equation}
Compute a set \(S\subseteq X\) by ordinary iteration from the empty set,
adjoining to the current \(S\)
\begin{itemize}
 \item every configuration that is not allowed;
 \item every \(t\in v\) for which there exist
       \(h\in T\) and \(w\in\mathcal V\setminus\{v\}\), with
       \(v\sim h\sim w\), such that
       \[
        \neg\exists t'\in w\setminus S\;\exists s\in h\,
                      [\Inc(t,s)\land\Inc(t',s)].
       \]
\end{itemize}
The old set \(S\) is retained at every update. Let \(S_\infty\) be the
fixed point. Every strict update adds an atom of \(X\), so stabilization
takes at most \(|X|+1\) applications of the update. Put
\begin{equation}
 \Phi_{\rm out}(\alpha,x)
 \quad\Longleftrightarrow\quad
 \operatorname{Complete}(\alpha,x)\ \land\
       \forall v\in\mathcal V\;(v\setminus S_\infty\ne\emptyset).
 \label{sen:output}
\end{equation}

\begin{lemma}[Exactness of the tree test]\label{sen:tree-test}
Suppose \(\operatorname{Check}(A)\) holds, \(\alpha\in P\), and
\(\operatorname{Complete}(\alpha,x)\) holds. Then
\(\Phi_{\rm out}(\alpha,x)\) is true if and only if there exist atoms
\(t_v\in v\), for every \(v\in\mathcal V\), and \(s_e\in e\), for every
\(e\in\mathcal E\), such that
\[
 \Inc(t_v,s_e)\quad\text{whenever }v\sim e,
 \qquad s_e\in x\quad\text{whenever }e\notin T.
\]
\end{lemma}
\begin{proof}
Suppose first that such atoms exist. Each \(t_v\) is allowed. Inductively,
none is deleted: across a tree edge from \(v\) to \(w\), the atom \(t_w\)
survives the previous round and shares \(s_e\) with \(t_v\).
Thus every block has a survivor.

Conversely, assume every block has a survivor at the fixed point.
All survivors are allowed, and every survivor has a surviving support
across each adjacent tree edge. Choose any survivor at the root of
\(T\). Proceed away from the root, choosing for each child a surviving
configuration supporting the one chosen at its parent. This defines
one \(t_v\) at every vertex because \(T\) is a finite tree. On each tree
edge, the two chosen configurations share an atom \(s_e\). Each
configuration has exactly one neighbor in each incident edge block by
\(\operatorname{Check}(A)\), so their selected edge atoms agree.
On a non-tree edge, allowedness forces both endpoints to use its unique
member of \(x\). These atoms give the required completion.
This argument does not require configurations within a block to be
determined uniquely by their incident neighbors.
\end{proof}

We can now specify the sentence:
\begin{equation}
 \begin{split}
 O(\alpha)
   &=\WSC^*xy.\,
       (s_{\rm step},s_{\rm choice},s_{\rm wit},\Phi_{\rm out}),\\
 \Phi
   &\equiv
       \{\emptyset:\alpha\in P,\ O(\alpha)\}\ne\emptyset.
 \end{split}
 \label{sen:sentence}
\end{equation}
The sole parameter of the WSC operator is the code \(\alpha\).
In particular, atom variables used to form the equivalence-class rows
are bound within those terms; none is an additional WSC parameter.
We give the fixed polynomial annotation below.

\subsection{Witnesses respect all previous states}

The guards make the sentence legal without any structural promise beyond
what it explicitly tests.

\begin{lemma}\label{sen:history}
For every input and every \(\alpha\in P\), each choice set reached from
\(x=\emptyset\) is a nonempty orbit under the automorphisms fixing
\(\alpha\) and every preceding state. The corresponding value of
\(s_{\rm wit}\) witnesses that orbit using automorphisms of the full
input. Every path reaches its first repeated state after at most
\(n+1\) transitions.
\end{lemma}
\begin{proof}
Inductively, the current state \(x\) consists of one selected atom from
each edge in an initial segment of the ordered set
\(\mathcal E\setminus T\). This holds at the empty state.
A real choice lies in the next unselected edge, so the step adds one new
atom and extends that initial segment. A dummy choice leaves \(x\)
unchanged and hence stabilizes. Thus all earlier states are subsets of
the current state. There are at most \(|\mathcal E|\leq n\) strict
steps, followed by a stabilizing step.

Fix an actual prefix \(b_0=\emptyset,b_1,\ldots,b_i=x\), and let
\[
 H_i=\Aut(A,\alpha,b_0,\ldots,b_i).
\]
Every element of \(M(\alpha,x)\) is a full-input automorphism fixing
\(\alpha\) and fixing all atoms in \(x\) individually. Since
\(b_j\subseteq x\) for \(j\leq i\), it fixes every earlier \(b_j\).
Consequently
\[
                         M(\alpha,x)\subseteq H_i.
\]
All terms defining \(D\) and \(M\) are equivariant in their displayed
arguments: they use only input relations, bounded set operations,
ordinals, and the indicated iterations. An automorphism fixing
\(\alpha\) and \(x\) therefore preserves \(D(\alpha,x)\), and it
conjugates the family \(M(\alpha,x)\) to itself.

If the first case of Equation~\eqref{sen:guarded-choice} applies, fix
\(s\in D(\alpha,x)\). Invariance under \(H_i\) gives
\(H_i s\subseteq D(\alpha,x)\), while the guard and
\(M(\alpha,x)\subseteq H_i\) give
\(D(\alpha,x)\subseteq H_i s\). Thus the choice set is exactly this
orbit, and the maps in \(M\) witness every ordered pair of its members.
If the second case applies, the choice set is \(\{\emptyset\}\),
and identity witnesses its only pair. It belongs to \(M\) in every
case. At a final state \(z\), the witnessing term evaluated on
\((\alpha,z,b_i)\) returns exactly \(M(\alpha,b_i)\), as required.
\end{proof}

The proof also shows directly why the output is independent of the
choices for a fixed \(\alpha\). Given two prefixes of the same length,
inductively map the first to the second by an automorphism fixing
\(\alpha\). At the next step, transport the first chosen member into the
second choice set by equivariance, and then use a witness fixing the
second prefix to send it to the second chosen member. Composing the
maps matches the extended prefixes. Stabilization is preserved by an
automorphism, so all terminal paths have the same length and all their
terminal states lie in one orbit. The formula \(\Phi_{\rm out}\) is
invariant on this orbit.

\subsection{One fixed polynomial for the expanded sentence}

All the notation in this section specifies finite syntax. In particular
the large integer \(B\) controls a fixed product. The raw witness
relations are enumerated by fixed products of atoms, and the WSC
semantics handles choice paths individually. We now give a bound that
also covers temporary relation graphs and objects rejected by the filter.

Fix a direct, capture-avoiding expansion of all the preceding macros
using the tuple convention already specified. Let \(H_0\) be its
syntax height; the number of symbols may be used instead as an upper
bound. Define the fixed integer sequence and polynomial
\begin{equation}
 d_0=1000,\qquad d_{j+1}=1000(d_j+1)^3,\qquad
 p(n)=(n+2)^{d_{H_0+1}}.
 \label{sen:polynomial}
\end{equation}
Every ordinary iteration and the WSC operator receive this same
polynomial annotation.

\begin{lemma}[Resource bound]\label{sen:resources}
For every finite input, the expansion of
Equation~\eqref{sen:sentence}, annotated by
Equation~\eqref{sen:polynomial}, has the following properties.
Every constructed set, including temporary sets before any filtering,
has transitive size at most \(p(n)\). Every ordinary iteration
stabilizes within \(p(n)\) steps, and every WSC path to first
stabilization has at most \(p(n)\) state labels.
\end{lemma}
\begin{proof}
We first bound the states of the iterations independently of the
annotation. The row families have at most \(n\) members, each a subset
of \(A\). Every state of the iteration defining \(R_J\) is a subset of
\(\mathcal V^2\times I_n\), with the fixed triple coding from above.
It has polynomial hereditary size even on an input failing
\(\operatorname{Check}(A)\). An inflationary iteration on this domain
has at most \(n^2(n+1)\) strict updates. Every state in the deletion
iteration is a subset of \(X\). By Lemma~\ref{sen:history}, every WSC
state is a subset of \(Y\), and a path to first stabilization has at most
\(n+2\) state labels. The parameter \(\alpha\) is a fixed-length tuple
of blocks, and each supplied choice is an atom or the pure dummy.
With the fixed pair coding, all these state and supplied-argument
hereditary sizes, and all these iteration lengths, are bounded by
\((n+2)^{100}\).

For the remaining terms, use hereditary size
\(\rho(a)=1\) for an atom and \(\rho(a)=1+|\TC(a)|\) for a set.
An actual activation here means a call to a subexpression with arguments
supplied during evaluation of this expanded sentence before resource
cutoffs. We include every choice path and every retrospective witness
call.
We verify by induction on expanded syntax height that a subexpression
of height at most \(j\), at any actual activation whose free arguments
have hereditary sizes at most \((n+2)^d\), where \(d\geq1\), constructs
individual values of hereditary size at most
\[
                              (n+2)^{d_jd}.
\]
This induction includes the evaluations forming all candidates in
\(\mathcal L,\mathcal C\), and \(M\), whether or not those candidates
pass a subsequent test.

At height zero, variables and the constants \(\emptyset,\Atoms\) satisfy
the bound. The primitive functions \(\Pair,\Union,\Unique,\Card\)
satisfy the next-height bound: pairing adds only finitely many wrappers
to the two hereditary contents, union and unique extraction use their
arguments' contents, and \(\Card\) produces an ordinal no larger than
the cardinality of the input set, or zero for an atom.

For a comprehension of height \(j+1\), its range has at most
\((n+2)^{d_jd}\) members, and each such member has hereditary size
bounded by the same quantity. Thus the condition and body have
argument exponent at most \(d_jd\). By induction each body value has
hereditary size at most \((n+2)^{d_j^2d}\).
The union of their hereditary contents, together with the comprehension
itself, is bounded by
\((n+2)^{(d_j^2+d_j+2)d}\). The recurrence for \(d_{j+1}\) dominates
this exponent and the bounds for all temporary condition and body
evaluations. Atomic formulas and Boolean combinations introduce no
larger sets.

At an ordinary iteration, its extra state argument has the independent
bound \((n+2)^{100}\) proved above. Replacing the argument exponent \(d\)
by \(100d\) bounds the step subexpressions by
\((n+2)^{100d_jd}\). The returned state has the independent state bound;
the recurrence again dominates both estimates.
At the WSC operator the same argument applies to its step, choice,
witness, and output subexpressions, because the intermediate and final
states and the choices have the stated independent bound. In
particular it applies when the witness term is reevaluated at any
earlier state. The operator supplies a truth value; it does not collect
the branching family of paths as one HF object.

This proves the induction. Starting with the sentence's empty free
assignment and using one extra height in
Equation~\eqref{sen:polynomial} gives the asserted common bound.
The independent iteration bounds also lie below this polynomial.
All these estimates were obtained before imposing resource cutoffs.
Therefore the annotated and unbounded evaluations agree at every
actual call, and no cutoff is reached.
\end{proof}

\begin{proposition}[A total all-input definition]\label{sen:all-input}
The sentence \((\Phi,p)\) of
Equations~\eqref{sen:sentence} and~\eqref{sen:polynomial} is a fixed
\(\CPT+\WSC\) sentence over \(\tau\). It evaluates to either true or
false on every finite \(\tau\)-structure, never to \(\dagger\), and its
true-model class is isomorphism-invariant.

More explicitly, it is true precisely when \(P\ne\emptyset\) and there
exists \(\alpha\in P\) such that every terminal state of the guarded
selection is complete and extends to consistent configurations as in
Lemma~\ref{sen:tree-test}.
\end{proposition}
\begin{proof}
All macros use only the permitted finite syntax and bounded iterations.
The only WSC calls in the outer comprehension have parameters in \(P\).
Lemma~\ref{sen:history} verifies every required witness using
automorphisms of the entire input, including all three
\(\Zrel_\delta\) relations, and respecting every earlier state.
Lemma~\ref{sen:resources} excludes exceeded bounds. Hence neither source
of \(\dagger\) occurs. Every path reaches a fixed point, so the output
semantics applies to a nonempty family of terminal states.
Equation~\eqref{sen:output} and Lemma~\ref{sen:tree-test} give the stated
characterization.

An isomorphism transports the block families, \(P\), all bounded
constructions, and all choice paths to their counterparts. It also
preserves the output test. It may permute the parameter tuples, which
leaves the outer existential comprehension unchanged. Thus the
true-model class is invariant. If the input check fails or \(P\) is
empty, the comprehension has empty range and the sentence is false.
This includes the empty input. A connected graph with one vertex and
no edges is also covered: its repeated four-entry tuple is resolving,
the dummy stabilizes immediately, and the tree test asks only for a
configuration in its nonempty vertex block.
\end{proof}
